\subsection{Criterion for semisimplicity}

PROPOSITION 2. --- Let \(\varrho\) be a unitary representation of G in a complex Hilbert space E. Let \(\mathfrak{F}\) be a filter base on \(\mathcal{M}_{c}(\mathrm{G})\) converging to the measure \(\varepsilon_{e}\) for the topology of compact convergence. Suppose that there exists \(\mathrm{M} \in \mathfrak{F}\) such that \(\varrho(\nu)\) is a compact endomorphism of E for every \(\nu \in \mathrm{M}\) .

The unitary representation \(\varrho\) is then semisimple and every irreducible unitary representation of G has finite multiplicity in \(\varrho\) .

Let us first prove a lemma.

Lemma 2. --- Let \(\varrho_{1}\) be a nonzero unitary representation of G isomorphic to a subrepresentation of \(\varrho\) . There exists a measure \(\nu \in M\) such that \(\varrho_{1}(\check{\nu} * \nu)\) is a nonzero compact Hermitian endomorphism of the space of \(\varrho_{1}\) . In particular, there exists a nonzero real number \(\lambda\) such that the eigenspace of \(\varrho_{1}(\check{\nu} * \nu)\) relative to \(\lambda\) is nonzero.

We may assume that \(\varrho_{1}\) is a subrepresentation of \(\varrho\) . Let E \(_{1}\) be its space. For every measure \(\nu \in \mathrm{M}\) , the endomorphism \(\varrho_{1}(\nu)\) is deduced from \(\varrho(\nu)\) by passing to subspaces. The hypothesis and prop. 3 of III, p. 5 therefore imply that \(\varrho_{1}(\nu)\) is compact.

Since \(\mathfrak{F}\) converges to the measure \(\varepsilon_{e}\) in \(\mathcal{M}_{c}(\mathrm{G})\) endowed with the topology of compact convergence and the space \(\mathrm{E}_{1}\) is nonzero, there exists \(\nu\in\mathrm{M}\) such that \(v=\varrho_{1}(\nu)\) is a nonzero endomorphism of \(\mathrm{E}_{1}\) (cf. n°2 of V, p. 400). The endomorphism \(u=\varrho_{1}(\check{\nu}*\nu)\) is equal to \(v^{*}\circ v\) (lemma 1 of V, p. 401); it is therefore nonzero (EVT, V, p. 39, prop. 2), Hermitian and compact, since \(v\) is compact.

Since \(u\) is Hermitian and nonzero, its spectrum is not reduced to zero (example 1 of I, p. 110). Finally, since \(u\) is compact, every \(\lambda \in \mathrm{Sp}(u)\setminus\{0\}\) is an eigenvalue of \(u\) (prop. 2 of III, p. 83).

Let us now prove the proposition.

We will use prop. 7 of V, p. 391 to establish that \(\varrho\) is semisimple. Let \(\varrho_{1}\) be a nonzero subrepresentation of \(\varrho\) and \(\mathrm{E}_{1} \subset \mathrm{E}\) its space; we must show that \(\varrho_{1}\) contains an irreducible subrepresentation.

Let \(\nu \in \mathrm{M}\) such that \(u = \varrho_{1}(\check{\nu}*\nu)\) is a nonzero compact Hermitian endomorphism of the space of \(\varrho_{1}\) and let \(\lambda\) be nonzero such that the eigenspace F of \(u\) relative to \(\lambda\) is nonzero (lemma 2). The space F is finite-dimensional (prop. 5 of III, p. 90). There exists a subrepresentation \(\varrho_{2}\) of \(\varrho_{1}\) on a subspace \(\mathrm{E}_{2}\) of \(\mathrm{E}_{1}\) such that \(\mathrm{E}_{2} \cap \mathrm{F}\) is nonzero and of minimal dimension. Let \(x\) be a nonzero element of \(\mathrm{E}_{2} \cap \mathrm{F}\) and let \(\mathrm{E}_{3}\) be the subrepresentation of \(\varrho_{1}\) generated by \(x\) . We then have \(\mathrm{E}_{3} \subset \mathrm{E}_{2}\) , whence \(\mathrm{E}_{3} \cap \mathrm{F} = \mathrm{E}_{2} \cap \mathrm{F}\) by minimality of the dimension of \(\mathrm{E}_{2} \cap \mathrm{F}\) .

Let us show that the representation \(E_{3}\) is irreducible. Let \(E_{4} \subset E_{3}\) be a closed subspace stable under G. We have \(E_{2} \cap F = (E_{4} \cap F) \oplus (E_{4}^{\circ} \cap F)\) , where \(E_{4}^{\circ}\) denotes the orthogonal complement of \(E_{4}\) in \(E_{3}\) (indeed, if \(y \in E_{2} \cap F\) , let \(y_{4} \in E_{4}\) be its orthogonal projection onto \(E_{4}\) ; since \(E_{4}\) and \(E_{4}^{\circ}\) are stable under u, the vector \(u(y_{4})\) is the orthogonal projection of \(u(y) = \lambda y\) onto \(E_{4}\) , whence \(u(y_{4}) = \lambda y_{4}\) , which means that \(y_{4} \in E_{4} \cap F\) , and the assertion follows). The minimality of the dimension of \(E_{2} \cap F\) then implies that either \(E_{4} \cap F = E_{2} \cap F\) , or \(E_{4}^{\circ} \cap F = E_{2} \cap F\) . In the first case, we have \(x \in E_{4}\) , whence \(E_{4} = E_{3}\) , and in the second, we have \(x \in E_{4}^{\circ}\) , whence \(E_{4}^{\circ} = E_{3}\) and \(E_{4} = \{0\}\) . The representation \(E_{3}\) is therefore irreducible.

We conclude by prop. 7 of V, p. 391 that the representation \(\varrho\) is semisimple.

Let \(\pi\) be an irreducible unitary representation of G whose multiplicity in \(\varrho\) is nonzero; let \(E_{\pi}\) be its space. There exists a measure \(\nu \in M\) and a nonzero real number \(\lambda\) such that \(u = \pi (\check{\nu} * \nu)\) is a nonzero compact Hermitian endomorphism of \(E_{\pi}\) and the eigenspace of \(u\) relative to \(\lambda\) is nonzero (lemma 2). Let F be the eigenspace of \(v = \varrho (\check{\nu} * \nu)\) relative to \(\lambda\) . For every subrepresentation \(E_1\) of E isomorphic to \(\pi\) , the endomorphism of \(E_1\) deduced from \(v\) by passing to subspaces is identified with \(u\) , and the intersection of \(E_1\) and F is nonzero. Thus, the multiplicity of \(\pi\) in \(\varrho\) is less than the dimension of the space F, which is finite (prop. 5 of III, p. 90).

Example. --- Suppose that G is unimodular, for example that G is a real semisimple Lie group. Let \(\Gamma \subset G\) be a discrete subgroup such that the quotient \(X = \Gamma \backslash G\) is compact. The group G acts on X on the right by multiplication. We denote by \(\beta\) the counting measure on \(\Gamma\) and we set \(\widetilde{\mu} = \mu / \beta\) (INT, VII, p. 44, § 2, n° 2, def. 1); this is a bounded G-invariant measure on X.

For every \(f \in \mathcal{L}^2(\mathrm{X},\widetilde{\mu})\) and every \(g \in \mathrm{G}\) , we define the function \(\varrho(g)f \in \mathcal{L}^2(\mathrm{X},\widetilde{\mu})\) by \(\varrho(g)f(x) = f(x \cdot g)\) . The map \(\varrho(g)\) is a continuous linear map, which induces by passing to quotients a unitary map on \(\mathrm{L}^2(\mathrm{X},\widetilde{\mu})\) , still denoted \(\varrho(g)\) . The map \(\varrho\) is a unitary representation of G in \(\mathrm{L}^2(\mathrm{X},\widetilde{\mu})\) (INT, VII, p. 135, § 2, n° 5, prop. 8).

Let \(\varphi \in \mathcal{K}(\mathrm{G})\) . For all compact subsets \(\mathrm{L}_1\) and \(\mathrm{L}_2\) of \(\mathrm{G}\) , the intersection \(\mathrm{T}\) of \(\Gamma\) and of the compact set \(\mathrm{L}_1\mathrm{Supp}(\varphi)\mathrm{L}_2^{-1}\) is finite. For every \((g,h) \in \mathrm{L}_1 \times \mathrm{L}_2\) , the series \(\sum_{\gamma \in \Gamma} \varphi(g^{-1}\gamma h)\) coincides with the finite sum \(\sum_{\gamma \in \mathrm{T}} \varphi(g^{-1}\gamma h)\) . Since \(\mathrm{G}\) is locally compact, the sum of this series, denoted \(k_{\varphi}(g,h)\) , is a continuous function on \(\mathrm{G} \times \mathrm{G}\) .

Let \((g,h) \in G \times G\) and let \((\gamma, \eta) \in \Gamma \times \Gamma\) . We have \(k_{\varphi}(\gamma g, \eta h) = k_{\varphi}(g, h)\) , hence \(k_{\varphi}\) defines by passing to the quotient a continuous function on \(X \times X\) , which we denote \(\widetilde{k}_{\varphi}\) . We also have \(\widetilde{k}_{\varphi} \in \mathcal{L}^{2}(X \times X, \widetilde{\mu} \otimes \widetilde{\mu})\) , since the space X is assumed compact.

We denote by \(\dot{x}\) the image in X of an element \(x\) of G under the canonical projection. Let \(\varphi \in \mathcal{K}(\mathrm{G})\) . For \(f \in \mathcal{L}^2(\mathrm{X},\widetilde{\mu})\) and \(x \in \mathrm{G}\) , we have

\[
\begin{array}{c} \varrho (\varphi) f (\dot {x}) = \int_ {\mathrm{G}} \varphi (g)   (\varrho (g) f) (\dot {x}) d \mu (g) = \int_ {\mathrm{G}} \varphi (g) f (\dot {x} \cdot g) d \mu (g) \\ = \int_ {\mathrm{G}} \varphi (x ^ {- 1} y) f (\dot {y}) d \mu (y) = \int_ {\Gamma \backslash \mathrm{G}} \Bigl (\sum_ {\gamma \in \Gamma} \varphi (x ^ {- 1} \gamma y) \Bigr) f (\dot {y}) d \widetilde {\mu} (\dot {y}) \end{array}
\]

(INT, VII, p. 46, § 2, no. 3, prop. 5). Since \(\widetilde{k}_{\varphi}\) belongs to the space \(\mathcal{L}^2 (\mathrm{X}\times \mathrm{X},\widetilde{\mu}\otimes \widetilde{\mu})\) , it follows that the endomorphism \(\varrho (\varphi)\) of \(\mathrm{L}^2 (\mathrm{X},\widetilde{\mu})\) coincides with the Hilbert-Schmidt endomorphism with kernel \(\widetilde{k}_{\varphi}\) ; this endomorphism is therefore compact (cor. 1 of III, p. 33).

There exists a sequence \((\varphi_{n})_{n\in\mathbf{N}}\) of functions in \(\mathcal{K}_{+}(\mathrm{G})\) with integral 1 such that \(\varphi_{n}\cdot\mu\) converges to \(\varepsilon_{e}\) in \(\mathcal{M}_{c}(\mathrm{G})\) endowed with the topology of compact convergence (INT, VIII, p. 139, § 2, no. 7, cor. 2). Prop. 2 consequently implies that the representation \(\varrho\) is semisimple and that the multiplicities of the irreducible unitary representations of G in \(\varrho\) are finite.

The irreducible unitary representations of G whose multiplicity in \(\varrho\) is nonzero are called \(\Gamma\) -automorphic representations of the group G.

